\section{Introduction}
\label{sec:introduction}

Contact-rich manipulation is partially observable in a particularly consequential
way. Vision can estimate object pose before contact, but it cannot directly reveal
the pressure redistribution that precedes slip, the direction of an incipient
constraint, or whether an apparently stable object is actually supported by the
left hand, the right hand, a tool, or the environment. Tactile sensing reduces
this ambiguity at an individual interface. Bimanual manipulation then creates a
second problem: the controller must relate several local interfaces whose states
are coupled through shared rigid-body motion, compliance, friction, and changes
in contact topology.

The distinction is easiest to see in a handover. Before transfer, the giver bears
most of the load. During dual contact, the same object couples both hands, and a
safe release depends on the receiver's support rather than on giver touch alone.
After release, the cross-hand path disappears. A representation that concatenates
two tactile streams can, in principle, approximate this mapping, but it does not
state which interfaces interact, which entity mediates the interaction, or why
the mapping should transfer to a new grasp topology. Conversely, a graph can name
the interfaces, yet a graph whose edges are fixed or supplied by a simulator does
not solve contact inference at deployment.

Recent work makes this review timely. Large tactile pretraining systems learn
sensor-robust local features from video-based touch, robot play, and multimodal
alignment~\cite{guzey2023tdex,higuera2025sparsh,yang2024unitouch,feng2025anytouch,zhao2024t3}.
Spatially anchored and cross-modal models seek to place touch in a geometric or
temporal context~\cite{huang2025sata,heng2025vitacformer}. Relational models use
graphs for tactile arrays, learned dynamics, morphology, and contact
physics~\cite{funabashi2022distributed,yang2023tacgnn,ai2024robopack,battaglia2016interaction,sanchez2020gns}.
At the same time, bimanual policy learning has expanded from simulated suites to
human demonstrations and real systems~\cite{chen2022bidexhands,grannen2023stabilize,shaw2024bidex,jiang2024dexmimicgen,li2025maniptrans}.
The 2026 frontier further includes tactile-reactive bimanual control, whole-hand
benchmarks, morphology-aware tactile policies, contact-anchored control, and
tool-oriented contact fields~\cite{niu2026trex,yuan2026ftp1,huang2026htbench,cui2026contactanchored,ma2026semanticcontact}.

The two late-August surveillance updates sharpen this frontier. Hierarchical
tactile world models now separate contact, deformation, and slip prediction;
sequential visuo-tactile learning uses temporal alignment and training-time
privileged posteriors; and physics-aware sensing explicitly models object
properties or time-dependent viscoelastic relaxation
~\cite{xue2026hitacwam,xu2026vtmuse,liu2026vitacphys,wang2026relaxation}.
Meanwhile, structured cloth dynamics combine a mesh prior with a learned global
residual~\cite{wang2026meshpriordit}. These advances strengthen the relevant
alternatives to a contact graph: a valid comparison is no longer ``graph versus
concatenation'', but observation-matched structured priors, predictive temporal
models, and residual global models. None of these additions, however, directly
demonstrates observation-grounded cross-hand load topology.

The September--early-October catch-up changes the wording of the gap itself.
BiView-Touch provides direct offline human cross-hand tactile representation
evidence, while temporal tactile encoding with compliance provides a strong real
bimanual handover control baseline~\cite{liang2026biviewtouch,marra2026temporaltactilehandover}.
CADeT additionally shows observation-grounded transmission-mode inference and
physical probing in visually observed tissue manipulation~\cite{hu2026cadet}.
These are positive results at different links of the evidence chain. The remaining
question is their combination with independently validated robotic cross-hand
loads, matched structural controls, and topology-defined transfer, rather than
whether relational touch or inferred interaction modes exist at all.

These advances do not yet answer a shared representation question:

\begin{quote}
\emph{When does an explicit, observation-grounded model of contact topology and
load propagation improve bimanual perception or control over an equally capable
unstructured temporal model?}
\end{quote}

This formulation avoids two shortcuts. The first is to equate tactile input with
physical grounding. A policy may receive fingertip forces while depending on
object state, simulator contacts, or a future reference trajectory. The second is
to equate a graph architecture with a scientific contribution. Generic relational
learning is well established~\cite{battaglia2016interaction,wang2018nervenet,ying2021graphormer};
the contribution must lie in what the nodes and edges mean, how they can be
inferred, which physical quantities they predict, and which counterfactual tests
they survive.

\subsection{Review objectives and contributions}

This review integrates literatures that are often treated separately: local
tactile representation learning, structured contact and dynamics models,
bimanual manipulation and handover, cooperative tool use, and benchmark design.
Its contributions are fourfold.

\begin{enumerate}
  \item It introduces an orthogonal taxonomy that separates task topology,
  representation level, structural prior, supervision source, deployment
  observability, and evaluation evidence. This prevents a method from appearing
  physically grounded solely because it uses tactile inputs or graph layers.

  \item It distinguishes \emph{local contact representation} from
  \emph{cross-interface coordination representation}. The former encodes what a
  sensor sees; the latter must preserve which interfaces are coupled, how support
  is redistributed, and how this relation changes over time.

  \item It positions quasi-static handover and cooperative tool--workpiece tasks
  as complementary experimental topologies. Handover isolates support transfer;
  tool use tests multi-hop, role-asymmetric coupling without making an unsupported
  claim that the area is empty.

  \item It converts broad calls for ``better physical representations'' into a
  falsifiable evaluation programme: matched observation and parameter budgets,
  oracle, inferred, and random structure comparisons; contact-transition metrics;
  object-, session-, and contact-sequence-level splits; uncertainty calibration;
  and targeted interventions.
\end{enumerate}

The intended outcome is not a leaderboard. Sensor form factors, embodiments,
training data, task semantics, and success definitions vary too much for a valid
pooled effect estimate. Instead, the review identifies which claims can be made
from current evidence and which experiments are still needed.

\subsection{Terminology}

We use \emph{tactile observation} for a measured taxel array, tactile image,
deformation field, event stream, or calibrated interface force/torque signal.
Simple simulator contact flags and full object state are called \emph{privileged
contact variables} unless the deployment system measures them. \emph{Contact
topology} denotes the set of active interfaces between hands, objects, tools,
workpieces, and the environment. \emph{Load path} denotes a sequence of such
interfaces through which forces and moments are transmitted. \emph{Structured
representation} covers graphs, contact fields, object-centric frames, spatial
anchors, and other models that state relations beyond token order. These terms
describe information and assumptions, not a preferred neural architecture.
